\chapter{Blocked Nose}
\index{Blocked Nose}

\section*{Definition and Overview}
A blocked nose (nasal congestion or obstruction) is the sensation that the nasal passages are stuffy or that breathing through the nose is difficult. It occurs when the lining of the nose becomes swollen, when mucus accumulates, or when the nasal passages are narrowed by the structures of the nose itself. A blocked nose is very common, usually short-lived, and most often caused by a cold or allergy. Persistent or one-sided blockage, however, can indicate a structural problem or, rarely, a more serious condition that requires assessment.

\section*{Epidemiology}
Nearly everyone experiences a blocked nose at some point. It is a leading symptom of the common cold, of which adults typically have several each year, and of allergic rhinitis, which affects roughly one in five people. Nasal polyps and a deviated nasal septum are more common in adults and become more likely with age. In young children, recurrent one-sided blockage frequently results from a foreign body pushed up the nose, and overuse of decongestant sprays is a growing cause of chronic congestion in adults.

\section*{Aetiology and Pathophysiology}
The nose becomes blocked through three main mechanisms, which often combine:

\begin{itemize}
    \item \textbf{Swelling of the lining:} viral infection or allergy inflames the nasal membranes, narrowing the air passages
    \item \textbf{Excess mucus:} secretions build up and further obstruct airflow
    \item \textbf{Structural narrowing:} a deviated septum, enlarged turbinates, or nasal polyps physically block the passage
    \item \textbf{Common causes:} colds and other viruses; allergic and non-allergic (vasomotor) rhinitis; pregnancy and hormonal changes; nasal polyps; a deviated septum; irritants such as smoke and dust; and, in children, an object pushed up the nose
    \item \textbf{Medication overuse:} decongestant sprays used for more than a few days produce rebound congestion, a condition called rhinitis medicamentosa
\end{itemize}

\section*{Clinical Features}
A blocked nose is often accompanied by other nasal symptoms:

\begin{itemize}
    \item A stuffy feeling and difficulty breathing through one or both nostrils
    \item A runny nose or thick mucus discharge
    \item Sneezing and an itchy nose, typical of allergy
    \item A reduced sense of smell and taste
    \item A sensation of pressure or fullness in the face
    \item Snoring, mouth breathing, and a dry mouth on waking
    \item Headache or a feeling of congestion in the sinuses
    \item In children, a persistent one-sided, often smelly discharge that suggests a foreign body
\end{itemize}

\section*{Differential Diagnosis}
The many causes of a blocked nose must be distinguished from one another:

\begin{itemize}
    \item Viral cold or influenza: sudden onset with other cold symptoms such as sore throat and malaise
    \item Allergic rhinitis: itching, sneezing, and watery discharge related to triggers
    \item Non-allergic (vasomotor) rhinitis: triggered by weather, smoke, or strong smells
    \item Acute or chronic sinusitis: facial pain, pressure, and coloured discharge
    \item Nasal polyps: persistent blockage with reduced sense of smell
    \item Deviated septum or enlarged turbinates: often one-sided blockage
    \item Rhinitis medicamentosa: from overusing decongestant sprays
    \item Nasal foreign body in children: one-sided, often smelly discharge
    \item Rarely, nasal tumours: persistent one-sided blockage with bleeding
\end{itemize}

\section*{When to Seek Medical Care}
A blocked nose usually settles on its own within a week or two, but medical advice should be sought when:

\begin{itemize}
    \item It persists for more than a few weeks or is persistently one-sided
    \item There is a thick, coloured nasal discharge with facial pain or fever
    \item There is nosebleeding or persistently blood-stained discharge
    \item A child has a one-sided, smelly discharge, suggesting a foreign body
    \item It is interfering badly with sleep, work, or daily life
    \item Decongestant sprays have been needed for more than a few days
\end{itemize}

\textbf{Contact your local emergency services for:} difficulty breathing, swelling of the face or throat, or any sign of a severe allergic reaction.

\section*{Diagnosis and Investigations}
Most cases are diagnosed clinically, with further tests reserved for persistent or unusual symptoms:

\begin{itemize}
    \item \textbf{History:} duration, which side is affected, triggers, character of discharge, and medicines used
    \item \textbf{Examination:} inspection of the nose with a light, and with a small camera (nasal endoscopy) in many practices
    \item \textbf{Allergy testing:} skin prick or blood tests when allergic rhinitis is suspected
    \item \textbf{Imaging:} a CT scan of the sinuses for persistent blockage, chronic sinusitis, or before sinus surgery
    \item \textbf{Samples:} a swab of nasal discharge in selected cases of suspected infection
    \item \textbf{Referral:} to an ear, nose, and throat specialist when structural causes are suspected
\end{itemize}

\section*{Management}
Most blocked noses can be managed at home, and the treatment depends on the cause:

\begin{itemize}
    \item \textbf{Saline washes:} salt-water rinses and sprays keep the nose moist and clear mucus
    \item \textbf{Steam inhalation:} steam from a bowl or hot shower helps loosen secretions
    \item \textbf{Decongestant sprays and tablets:} effective in the short term, but sprays must not be used for more than a few days
    \item \textbf{Antihistamines:} for allergic rhinitis
    \item \textbf{Steroid nasal sprays:} the mainstay of treatment for allergy and polyps, used regularly over several weeks
    \item \textbf{Fluids and rest:} for colds
    \item \textbf{Avoiding triggers:} smoke, dust, and other irritants
    \item \textbf{Surgery:} septoplasty, turbinate reduction, or polyp removal for structural blockage that does not respond to medical treatment
\end{itemize}

\section*{Complications}
\begin{itemize}
    \item Sinusitis, when a blocked nose traps mucus and bacteria in the sinuses
    \item Middle-ear problems, especially in children, when the connection to the ear becomes blocked
    \item Sleep disturbance, snoring, and poor-quality sleep
    \item A reduced sense of smell
    \item Dependence on decongestant sprays (rhinitis medicamentosa)
    \item A missed diagnosis of structural or, rarely, more serious disease with persistent one-sided blockage
\end{itemize}

\section*{Prognosis and Outcomes}
A blocked nose from a cold usually improves within one to two weeks. Allergic rhinitis is a chronic condition but is usually well controlled with regular nasal sprays and trigger avoidance. Nasal polyps and septal problems respond well to surgery when medical treatment is insufficient, and most people achieve good nasal breathing with the right combination of treatments. Rebound congestion settles once decongestant sprays are stopped, although this may take some time and patience.

\section*{Prevention and Self-Care}
\begin{itemize}
    \item Wash hands regularly to reduce the frequency of colds
    \item Avoid known allergy triggers, such as pollen, dust mites, and pet dander
    \item Use a saline nasal wash regularly if you are prone to congestion
    \item Limit decongestant spray use to a maximum of a few days
    \item Avoid smoking and smoky environments
    \item Use a humidifier or steam inhalation if the air is dry
    \item Sleep with the head elevated if congestion is worse at night
    \item Keep small objects away from young children
\end{itemize}

\section*{Sources and Further Reading}
The following sources provide reliable, up-to-date information for patients and students of medicine.

\begin{itemize}
    \item NHS. \textit{Overview of a blocked or stuffy nose and how to treat it.} https://www.nhs.uk/conditions/blocked-or-runny-nose/
    \item Mayo Clinic. \textit{Nasal congestion -- causes and home remedies.} https://www.mayoclinic.org/
    \item MSD Manual. \textit{Nasal congestion and discharge -- professional edition.} https://www.msdmanuals.com/
    \item MedlinePlus. \textit{Nasal congestion -- patient information.} https://medlineplus.gov/
    \item NICE. \textit{Clinical knowledge summaries on allergic rhinitis and nasal congestion.} https://cks.nice.org.uk/
\end{itemize}
